\documentclass[11pt]{article}
\usepackage[margin=1in]{geometry}
\usepackage{amsmath,amssymb,amsthm}
\usepackage[hidelinks]{hyperref}

\newtheorem{theorem}{Theorem}
\newtheorem{proposition}[theorem]{Proposition}
\newtheorem{lemma}[theorem]{Lemma}
\theoremstyle{definition}
\newtheorem{remark}[theorem]{Remark}

\newcommand{\R}{\mathbb{R}}
\newcommand{\norm}[1]{\left\|#1\right\|}

\title{Bernoulli Expansion Proofs: Corrected Status}
\author{Plan 1 / Sharp stability for Faber--Krahn}
\date{2026-05-11}

\begin{document}
\maketitle

\begin{abstract}
This note corrects the overclaims in the previous Bernoulli expansion proof.
The Bernoulli spectral idea remains
promising, but the earlier writeup did not prove the required
Schauder/tame remainder estimates. In particular, it incorrectly evaluated the
interior harmonic first variation outside \(B_1\). The note below records the
correct fixed-domain formulation and the conditional spectral closure that
would follow from the missing tame estimates.
\end{abstract}

\section{What Was Wrong}

The previous proof defined
\[
u_1(r,\theta)=\frac1N\sum_{k\ge0}r^k g_k(\theta)
\]
as the harmonic extension of \(g/N\) in \(B_1\), then evaluated it at
\(r=1+g(\theta)\). For a general \(g\in C^{2,\gamma}\), this series is only an
interior harmonic expansion. It has no controlled extension to \(r>1\). Thus
the claimed estimates for
\[
u_1((1+g)\theta),\qquad \partial_r u_1(1+g,\theta)
\]
are not valid.

Consequently, the claimed quadratic Schauder remainder
\[
\|u_g-u_{B_1}-u_1\|_{C^{2,\gamma}}\lesssim \|g\|_{C^{2,\gamma}}^2
\]
on \(U_g\), and the derived tame estimate
\[
\|R_q(g)\|_{L^2}
\lesssim
\|g\|_{C^{2,\gamma}}\|g\|_{L^2},
\]
were not proved.

\section{Correct Fixed-Domain Formulation}

Let
\[
\Phi_g(r\theta)=(r+\chi(r)g(\theta))\theta,\qquad
\widetilde u_g=u_g\circ\Phi_g.
\]
Then \(\widetilde u_g=0\) on \(\partial B_1\), and
\[
-\operatorname{div}(A_g\nabla\widetilde u_g)=J_g
\quad\hbox{in }B_1,
\qquad
A_g=J_gD\Phi_g^{-1}D\Phi_g^{-T}.
\]
This is the right framework. The first variation of the shape derivative
\(u'\) still satisfies
\[
\Delta u'=0\quad\hbox{in }B_1,\qquad
u'|_{\partial B_1}=\frac{g}{N}.
\]
But the pullback derivative is not simply \(u'\circ\Phi_g\) evaluated outside
the ball. Boundary-gradient expansions must be derived on the fixed boundary
from the pulled-back equation and the normal transformation formula.

\section{The Useful Linearization}

The first-order Bernoulli map is still expected to be
\[
q_g
=
\frac1N-\frac1N(\mathcal L g)+\hbox{higher order},
\qquad
\mathcal L g_k=(k-1)g_k.
\]
This part is consistent with standard shape calculus:
\begin{itemize}
\item \(u'\) has boundary data \(g/N\);
\item the Dirichlet-to-Neumann map on the unit ball sends \(g_k\) to \(k g_k\);
\item evaluating the normal derivative on the moved boundary contributes the
additional \(-g_k\), giving \(k-1\).
\end{itemize}
The \(\alpha\)-variation bracket also still has the linearization
\[
\Psi_g=1+(I-\Pi_1)g+\hbox{higher order},
\]
where \(\Pi_1\) is the degree-one spherical-harmonic projection.

\section{Conditional Spectral Closure}

\begin{proposition}[Conditional Bernoulli spectral closure]
Assume that in the small \(C^{2,\gamma}\) graph class,
\[
q_g=\frac1N-\frac1N\mathcal L g+R_q(g),
\qquad
\|R_q(g)\|_{L^2}\le C\|g\|_{C^{2,\gamma}}\|g\|_{L^2},
\]
and
\[
\Psi_g=1+(I-\Pi_1)g+R_\alpha(g),
\qquad
\|R_\alpha(g)\|_{L^2}\le C\|g\|_{C^1}\|g\|_{L^2}.
\]
Then, for selected minimizers satisfying
\[
q_g^2=\Lambda+a_\sigma\Psi_g,\qquad |a_\sigma|\le C\sigma,
\]
small \(\sigma\) and small \(\|g\|_{C^{2,\gamma}}\) force \(g=0\), after
volume and barycenter normalization.
\end{proposition}

\begin{proof}
Squaring the assumed Bernoulli expansion and subtracting the constant part gives
\[
-\frac{2}{N^2}\mathcal Lg+R(g)
=a_\sigma(I-\Pi_1)g+a_\sigma R_\alpha(g),
\]
with
\[
\|R(g)\|_{L^2}\le C\|g\|_{C^{2,\gamma}}\|g\|_{L^2}.
\]
Projecting to modes \(k\ge2\) gives
\[
\mathcal Lg_{\ge2}
=O(\sigma)g_{\ge2}
+O(\|g\|_{C^{2,\gamma}})\,g .
\]
On this subspace \(\|\mathcal Lh\|_{L^2}\ge\|h\|_{L^2}\), so
\[
\|g_{\ge2}\|_{L^2}
\le C(\sigma+\|g\|_{C^{2,\gamma}})\|g\|_{L^2}.
\]
Volume and barycenter normalization give
\[
\|g_0\|_{L^2}+\|g_1\|_{L^2}\le C\|g\|_{L^2}^2.
\]
For small enough \(\sigma\), \(\|g\|_{C^{2,\gamma}}\), and \(\|g\|_{L^2}\), this
forces \(g=0\).
\end{proof}

\section{What Remains Genuinely Open}

To make the Bernoulli spectral route publication-ready, one must prove the
tame Bernoulli expansion without evaluating the harmonic first variation
outside \(B_1\). A viable approach is:
\begin{enumerate}
\item work entirely on \(B_1\) with the pullback equation;
\item differentiate the coefficient map \(g\mapsto(A_g,J_g,\nu_g)\);
\item prove a second-order fixed-domain remainder estimate for
\(\widetilde u_g\) in a norm strong enough to control the boundary normal
derivative;
\item show the boundary-gradient remainder is tame in the mixed form
\(C^{2,\gamma}\times L^2\), not merely \(C^{2,\gamma}\)-quadratic.
\end{enumerate}
Until that is done, the selection route should still be considered closed by
BDV's nearly spherical second variation, not by the Bernoulli spectral closure.

\end{document}
